\chapter{绪论}

\section{研究背景及意义}

近三十年来，实时计算机图形学领域长期由光栅化（Rasterization）技术主导。光栅化凭借其卓越的并行计算效率，通过将三维几何顶点投影至二维屏幕空间并进行片元插值，能够在有限的硬件算力下实现高分辨率与高帧率的图像渲染。然而，光栅化本质上属于基于局部几何信息的启发式算法，难以精确模拟物理世界中光线在三维空间内的全局传输与多次反弹路径。因此，在处理全局光照（Global Illumination, GI）、基于物理的软阴影（Soft Shadows）以及多重粗糙镜面反射（Glossy Reflections）等复杂光学现象时，传统光栅化管线通常只能依赖预计算或屏幕空间经验近似算法。例如，使用屏幕空间环境光遮蔽（SSAO）来模拟暗角，使用级联阴影贴图（CSM）来计算硬阴影，以及使用屏幕空间反射（SSR）来模拟镜面效果。这些近似方案在复杂场景中极易引入明显的视觉伪影（Visual Artifacts），如 SSR 导致的屏幕外信息缺失与反射断裂。此外，随着现代数字娱乐产业对渲染真实感需求的日益提升，堆叠各类屏幕空间特效往往会导致渲染管线复杂度剧增，系统维护成本急剧上升。

随着图形硬件算力的持续突破以及 Vulkan Ray Tracing、DirectX Raytracing (DXR) 等现代图形 API 扩展的标准化，硬件加速的实时光线追踪（Hardware-Accelerated Ray Tracing）技术已正式步入实时渲染的舞台。该技术通过在硬件层面加速边界体积层次结构（Bounding Volume Hierarchy, BVH）的遍历与射线-三角形求交测试，能够精确求解渲染方程（Rendering Equation），从而生成具有照片级真实感（Photorealism）的图像。然而，受限于当前图形处理器（GPU）的显存带宽与核心算力瓶颈，在复杂的高分辨率实时交互场景中，若直接采用与离线影视渲染相同的纯路径追踪（Path Tracing）算法，每像素需要发射数千条光线，这与实时渲染严格的 $16.6$ 毫秒（60 FPS）帧时间预算存在巨大的性能鸿沟。

面对这一性能与画质的尖锐矛盾，混合渲染（Hybrid Rendering）架构应运而生。混合渲染的核心思想是“扬长避短”：利用光栅化极高的可见性判定效率生成包含场景几何与材质属性的几何缓冲（G-Buffer），随后利用硬件光线追踪在 G-Buffer 的表面属性基础上发射极少量（通常为 1 SPP，即每像素 1 条射线）光线，专门用于求解光栅化难以处理的阴影、反射和间接光照信号。然而，1 SPP 的极低采样率必然会导致渲染结果呈现出剧烈的蒙特卡洛随机噪声。因此，必须通过先进的时空降噪算法（Spatiotemporal Denoising）与时间抗锯齿技术（Temporal Anti-Aliasing, TAA）来修复低采样率带来的信噪比缺失。本文正是在此背景下，基于 Vulkan API 深入探索并实现了一套具备工业级降噪与抗锯齿能力的实时光线追踪混合渲染系统。

\section{国内外研究现状与发展}

实时渲染技术的发展史，本质上是对物理世界光照规律不断逼近的演进史。当前国内外的研究现状主要集中在以下三个技术脉络：

\subsection{实时渲染框架的演进}
早期的实时渲染引擎主要采用前向渲染（Forward Rendering）管线。前向渲染在处理每一个几何图元时，都会立即遍历场景中的所有光源进行光照计算。这种几何绘制与光照计算高度耦合的模式，导致其在面对多光源、高复杂度场景时，由于大量不可见片元被重复计算（Overdraw），性能会急剧下降。

为了解决多光源的性能瓶颈，延迟渲染（Deferred Rendering）被提出并广泛应用于现代 3D 引擎（如 Unreal Engine、Unity）中。延迟渲染通过引入几何缓冲（G-Buffer），将几何处理阶段与光照计算阶段完全分离。在几何阶段，系统仅将可见图元的反照率（Albedo）、法线（Normal）、材质粗糙度（Roughness）和深度（Depth）等属性写入 G-Buffer；在光照阶段，则直接利用 G-Buffer 中的物理属性进行全屏的光照合成。这种架构将光照计算的复杂度从“光源数量 $\times$ 几何图元数量”降低为“光源数量 $\times$ 屏幕像素数量”，极大提升了复杂光照环境下的渲染效率。然而，延迟渲染也带来了显存带宽开销大、难以处理半透明物体以及与传统多重采样抗锯齿（MSAA）不兼容等局限性。

\subsection{光线追踪技术的发展}
光线追踪技术最早由 Turner Whitted 于 1980 年提出（Whitted-style Ray Tracing），奠定了利用递归射线追踪实现镜面反射与折射的理论基础。1986 年，Kajiya 提出了著名的渲染方程（Rendering Equation），并将蒙特卡洛积分引入光线追踪，形成了路径追踪（Path Tracing）算法，首次在数学上统一了全局光照的求解模型。

长期以来，受限于庞大的计算量，该技术仅应用于影视特效等离线渲染领域。直到 2018 年，NVIDIA 发布了包含 RT Core 专用硬件加速单元的 RTX 系列显卡，Khronos Group 随后推出了 \texttt{VK\_KHR\_ray\_tracing} 扩展集，实时光线追踪才在底层硬件和 API 层面成为可能。目前，国内外顶尖工业界的主流做法均是结合光栅化与光追的混合架构，但如何在极低的光线预算下最大化光追画质，仍是学术界和工业界共同探索的热点方向。

\subsection{实时降噪与抗锯齿技术的发展}
在混合渲染管线中，降噪（Denoising）与抗锯齿（Anti-Aliasing）是决定最终画质的核心环节。
\begin{itemize}
    \item \textbf{抗锯齿技术}：由于延迟渲染无法高效使用硬件 MSAA，早期工业界多采用快速近似抗锯齿（FXAA）或次像素形态抗锯齿（SMAA）等纯空间后处理算法，但这些算法无法解决高频几何边缘的时间闪烁问题。近年来，时间抗锯齿（Temporal Anti-Aliasing, TAA）逐渐成为主流。TAA 通过在投影矩阵中引入 Halton 亚像素抖动（Sub-pixel Jitter），结合运动矢量（Motion Vector）将多帧的历史采样结果进行指数滑动平均（EMA），以极低的性能开销实现了媲美超采样（SSAA）的抗锯齿质量。
    \item \textbf{降噪技术}：针对 1 SPP 极低样本率下的蒙特卡洛噪声，传统的空间降噪算法（如双边滤波或 Edge-avoiding \`{A}-Trous 小波滤波）在强力抹平噪声的同时，往往会严重模糊几何边界与纹理细节。2017 年，Schied 等人提出了时空方差导向滤波（SVGF, Spatiotemporal Variance-Guided Filtering）算法。该算法开创性地将时间域的历史累积与空间域的 \`{A}-Trous 滤波相结合，通过估算局部时间方差（Temporal Variance）来动态调整空间滤波的边缘停止权重（Edge-stopping weights）。SVGF 成功实现了在保留高频几何细节的同时抹平低频噪声，目前已成为业界处理实时光追信号的标准基线算法。
\end{itemize}

\section{本文主要研究内容与贡献}

本文针对现代实时渲染中高画质与高性能难以兼顾的问题，基于 Vulkan 图形 API 设计并实现了一套完整的实时光线追踪混合渲染系统。针对工业界开发中的痛点，本文在系统架构、混合管线设计以及时空信号处理上进行了深度优化。主要工作与贡献如下：

\begin{enumerate}
    \item \textbf{设计并实现了基于 Render Graph 的渲染图架构}：针对 Vulkan 显式 API 中显存屏障（Memory Barrier）、图像布局转换（Image Layout Transition）和生命周期管理极其繁琐且易出错的痛点，实现了一套自动化的渲染图调度系统。该系统能够根据资源的读写约束和有向无环图（DAG）拓扑排序，在运行时自动推导管线的同步状态，极大提升了混合渲染管线的扩展性与资源复用率。
    
    \item \textbf{构建了高效的混合物理光照渲染管线}：结合延迟渲染框架，在 G-Buffer 的基础上实现了基于物理的渲染（PBR）光照计算。针对不同的光照特征灵活采用了底层光追技术：利用 Ray Query 结合几何偏导数（Derivative Face Normal）实现无自相交伪影的精确软硬阴影；利用 RT Pipeline（Raygen, Miss, ClosestHit）实现硬件加速的粗糙镜面反射与漫反射全局光照（GI）。
    
    \item \textbf{实现了基于 SVGF 的工业级时空降噪系统}：针对 1 SPP 光追信号的严重噪声问题，实现并优化了 SVGF 降噪数据流。本文深入解决了 SVGF 算法在实际工程中遇到的多项隐蔽缺陷：包括修正 Reversed-Z（反转 Z 缓冲）非线性深度导致的边缘判定失效问题，以及提出在光追着色器中解耦反照率（Demodulate Albedo），从而避免空间滤波对纹理细节的错误涂抹。
    
    \item \textbf{集成了高精度的 TAA 时间抗锯齿技术}：基于 Halton 序列实现了亚像素抖动系统。为解决 TAA 重投影过程中的历史错位与屏幕物理震动问题，本文在着色器底层推导并实现了严谨的抖动补偿公式（Jitter Compensation），并通过 AABB 颜色方差裁剪（Variance Clipping）有效消除了动态场景下的鬼影（Ghosting）现象，最终输出了极其平滑、稳定的照片级渲染画面。
\end{enumerate}

\section{论文结构安排}

本文的组织结构如下：
\begin{itemize}
    \item \textbf{第一章 绪论}：阐述课题的研究背景、意义及国内外发展现状，并总结本文的主要研究内容与贡献。
    \item \textbf{第二章 实时混合渲染理论基础}：详细介绍基于物理的渲染（PBR）材质模型、Vulkan API 的核心机制与显式同步理论，以及硬件光线追踪加速结构（BVH）的基本原理。
    \item \textbf{第三章 基于 Render Graph 的渲染器架构设计}：详细论述渲染图的节点抽象、虚拟资源池管理、以及 Vulkan 状态屏障的自动推导算法与乒乓轮换（Ping-Pong）机制的工程实现。
    \item \textbf{第四章 混合光照管线的核心设计与实现}：重点分析 G-Buffer 延迟管线构建、微表面 BRDF 物理光照合成逻辑，以及 Ray Query 与 RT Pipeline 在具体光影效果中的应用与优化。
    \item \textbf{第五章 基于时空域的抗锯齿与降噪算法实现}：深入剖析蒙特卡洛噪声与锯齿的形成原理。详细推导 TAA 的亚像素重投影数学模型，并全面讲解 SVGF 算法中历史累积、方差估算与 \`{A}-Trous 空域滤波的实现细节及工业级防坑指南。
    \item \textbf{第六章 渲染效果展示与性能分析}：通过丰富的对比实验图展示混合渲染器的最终画面质量，对比光栅化与光线追踪的视觉差异，并通过 GPU 耗时统计验证系统在不同负载下的实时性能。
    \item \textbf{第七章 总结与展望}：对全文的核心工作进行总结，客观分析当前系统的局限性，并对引擎未来可能的优化方向（如 ReSTIR 等新兴算法）进行展望。
\end{itemize}